\label{sec:history}

\subsection{1790: The First Apportionment}

The first U.S.\ congressional apportionment followed the 1790 Census.
The Apportionment Act of 1792 used Jefferson's method with a divisor
of 33{,}000, producing the following seat counts
\citep{schmeckebier1941congressional}:
Virginia 19 seats, Massachusetts 14, Pennsylvania 13, New York 10,
and so on.
Jefferson's role in recommending his method placed him at the intersection
of political self-interest (Virginia, the largest state, benefited from
floor rounding) and mathematical analysis.
President Washington exercised his first veto against an earlier
apportionment bill on constitutional grounds, and Jefferson's method
was adopted as the compromise \citep{balinski1982fair}.

\subsection{1790--1840: Six Apportionments}

Jefferson's method was used without revision for the apportionments
following the 1790, 1800, 1810, 1820, 1830, and 1840 censuses.
The divisor was adjusted each decade to achieve the desired House size.
Throughout this period, the method consistently overrepresented
Virginia and later New York (the largest states at different times)
relative to their exact proportional quotas.

Key instances of bias:
\begin{itemize}
  \item \textbf{1800 Census}: Virginia received 19 seats with an
    exact quota of 18.33 --- overrepresentation of 0.67 seats.
  \item \textbf{1820 Census}: New York received 34 seats with an
    exact quota of 32.5 --- overrepresentation of 1.5 seats.
  \item \textbf{1830 Census}: New York received 40 seats with an
    exact quota of approximately 38.59 --- overrepresentation of
    more than 1 seat, violating the upper quota rule.
\end{itemize}

\subsection{The 1832 Controversy and Abandonment}

The 1832 controversy was decisive.
Webster, then a senator from Massachusetts, introduced a bill replacing
Jefferson's method with his major fractions method.
His argument was mathematical: Jefferson's method allowed a state to
receive more seats than its upper quota, which was an arithmetic
error, not a policy choice.
The Senate passed Webster's bill; the House rejected it.
The 1840 apportionment still used Jefferson's method.

After the 1840 Census, the 1842 Apportionment Act replaced Jefferson
with Webster's method.
Jefferson was never subsequently used for federal apportionment in
the United States.

\subsection{D'Hondt's Independent Invention}

Victor d'Hondt, a Belgian lawyer and mathematician, independently
derived the same method in 1882 for the purpose of allocating seats
among parties in the Belgian Parliament \citep{dhondt1882systeme}.
His paper, written in French and published in Belgium,
was unknown to U.S.\ apportionment scholars for decades.
The equivalence between Jefferson and d'Hondt was noted in the
comparative electoral systems literature of the 1960s and 1970s
and fully formalised by \citet{balinski1982fair}.

The historical independence of the two derivations underscores the
mathematical naturalness of the method: both Jefferson (concerned with
state apportionment) and d'Hondt (concerned with party-list PR)
arrived at the same priority formula $P_i(s) = p_i/s$
as the most natural divisor method.
The floor rounding interpretation (Jefferson/d'Hondt = floor-rounding
divisor method) was not explicit in either original paper but
was established by later mathematical analysis \citep{balinski1982fair}.
